\section*{Chapitre \ref{ch:ml:text-from-bag-of-words-to-embeddings} --- Texte~: du sac de mots aux plongements}

\begin{solution}{exo:ml:text-from-bag-of-words-to-embeddings:1}
Jetons~: \texttt{<co>}, did, not, cut, guidance. Bigrammes~: \texttt{<co> did}, did not, not cut, cut guidance. «~Not cut~»
porte la négation~; le mot «~cut~» seul se lit comme une mauvaise nouvelle.
\end{solution}

\begin{solution}{exo:ml:text-from-bag-of-words-to-embeddings:2}
$\log(20\,000/3\,000) = 1.90$ et $\log(20\,000/200) = 4.61$~: «~halves~» pèse environ 2.4 fois autant. Les mots rares sont
plus informatifs sur le document dont il s'agit, ce qui n'est pas la même chose que plus informatifs sur le rendement.
\end{solution}

\begin{solution}{exo:ml:text-from-bag-of-words-to-embeddings:3}
$(1 - 1)/6 = 0$ pour la phrase composée («~beat~» positif, «~cut~» négatif), effet injecté $-1.2~\%$~; $1/4 = 0.25$
pour la récompense («~wins~» positif), effet injecté 0.
\end{solution}

\begin{solution}{exo:ml:text-from-bag-of-words-to-embeddings:4}
Ils apparaissent à la même place, avec les mêmes voisins (une société, puis «~estimates~», «~forecasts~» ou «~expectations~»), et
la cooccurrence est tout ce que voit le plongement~: cosinus 0.95. Seul un signal qui les distingue les sépare~: les rendements (un
\omterm{def:ml:representation-learning:ssl}{ajustement fin} supervisé), ou des contextes qui diffèrent («~beat \dots\ shares rise~»), ce qu'un corpus plus vaste fournit en partie.
\end{solution}

\begin{solution}{exo:ml:text-from-bag-of-words-to-embeddings:5}
Ce n'est pas un bug ici~: dans ce corpus, «~to be~» n'apparaît que dans les titres sur des rachats, donc il porte leur rendement. Il devient un bug
lorsque la formulation qui le rendait informatif change (le style d'une nouvelle agence, un nouveau gabarit), ou lorsque le mot désigne
quelque chose hors du texte, comme la couverture d'un média pour un type d'entreprise~; un mot repéré par le filtre doit être relu, et
la liste vérifiée pour sa stabilité d'un pli à l'autre.
\end{solution}

\begin{solution}{exo:ml:text-from-bag-of-words-to-embeddings:6}
Le fichier de symboles d'aujourd'hui ne fait correspondre les anciens tickers à rien et les tickers réutilisés au titulaire actuel~; les erreurs se concentrent dans le passé,
dans les sociétés qui ont changé ou disparu, précisément là où vivent les biais de survie et les effets de rachat. Un contrôle manuel
des liens récents, ou des liens qui ont abouti, mesure la mauvaise population. Ici, les alias d'aujourd'hui donnent une précision de 97.5~\% parmi
les mentions liées, mais seulement 87.5~\% de mentions correctement liées au total, et une précision de 91.5~\% sur les tickers. Vérifier un échantillon stratifié par année et par
type de mention, y compris les mentions non liées, en le comparant aux alias en vigueur à chaque date.
\end{solution}

\begin{solution}{exo:ml:text-from-bag-of-words-to-embeddings:7}
Sans masquage, l'IC du modèle TF-IDF passe de 0.273 à 0.266~: les mots de société reçoivent de petits poids ajustés au bruit, puisque aucun
effet de société n'a été injecté. Le filtre supervisé ne retient aucun mot de société et obtient 0.238 dans les deux cas, parce que son
seuil à trois erreurs types les rejette. Sur des données réelles, les mots de société porteraient les rendements de chaque société sur la période d'entraînement
(une bonne année, un rachat), une fuite de performance passée déguisée en texte, que le masquage empêche.
\end{solution}

\begin{solution}{exo:ml:text-from-bag-of-words-to-embeddings:8}
Pour un modèle linéaire d'ordonnée à l'origine $b$ et de poids des mots $w$, la phrase composée «~$A$ but $C$~» obtient le score $b + w_A + w_{\text{but}} + w_C
= f(A) + f(C) + (w_{\text{but}} - b)$, où $f(A)$ et $f(C)$ sont les scores des titres à une seule proposition avec les mêmes
mots. Pour que $f(\text{beat}) = 1$, $f(\text{cut}) = -1.5$ et $f(\text{«~beat but cut~»}) = -1.2$, il faut $w_{\text{but}} - b =
-0.7$~; pour que $f(\text{missed}) = -1$, $f(\text{raised}) = 1.5$ et $f(\text{«~missed but raised~»}) = 1.2$, il faut $+0.7$. Les
deux conditions se contredisent. Des bigrammes comme «~but cut~» et «~but raised~», ou une interaction entre la proposition
après «~but~» et un indicateur de guidance, peuvent satisfaire les deux~: la seconde proposition domine.
\end{solution}

\begin{solution}{pb:ml:text-from-bag-of-words-to-embeddings:1}
\textbf{Partie I.}
\begin{enumerate}
\item Pour qu'aucun modèle n'apprenne l'identité d'une société (ses rendements passés dans l'échantillon) comme signal de texte, et pour qu'un seul modèle
serve toutes les sociétés.
\item Il reçoit le poids $\log(N/N) = 0$~: il ne peut pas distinguer les documents.
\item De 123 mots à 473 variables, soit presque quatre fois plus~; les phrases composées qu'elles peuvent capturer représentent 9~\% des titres,
et la variance ajoutée a coûté plus que l'ajustement n'a gagné.
\item L'effet injecté explique tout ce que le texte peut expliquer~: un IC d'environ 0.3, parce que le bruit de la réaction (3~\%) est grand face
aux effets (jusqu'à 2~\%). Tout score au-dessus sur un \omterm{def:ml:why-financial-machine-learning-is-different:sets}{ensemble de test} relève de la chance ou d'une fuite.
\end{enumerate}
\textbf{Partie II.}
\begin{enumerate}[resume]
\item Règles d'événements 0.292, mots TF-IDF 0.273, TF-IDF avec bigrammes 0.257, mots supervisés 0.238, dictionnaire 0.224,
plongements 0.094.
\item Le modèle TF-IDF sur les mots dépasse le dictionnaire entre 1\,000 et 3\,000 titres étiquetés~; la liste supervisée et
le modèle à bigrammes seulement au-delà de 3\,000.
\item Elles ont été écrites à partir des gabarits propres au corpus et lisent correctement les 20\,000 titres~; elles échouent sur une formulation
pour laquelle elles n'ont pas été écrites («~surpasses consensus~» ne renvoie rien). Il faut les mesurer sur des textes rédigés après les règles.
\item Les plongements appris à partir de la cooccurrence rapprochent des contraires («~beat~» et «~missed~» à 0.95), si bien que la moyenne des
vecteurs d'un titre distingue à peine les bonnes des mauvaises nouvelles.
\end{enumerate}
\textbf{Partie III.}
\begin{enumerate}[resume]
\item 99.5~\%~; ses seules erreurs concernent le nom court Apex sans mot de secteur, où l'Apex le plus couvert est retenu.
\item 87.5~\%. Les anciens noms et tickers ne renvoient à rien (89.7~\% des mentions liées), et les tickers réutilisés renvoient au nouveau
titulaire (précision de 91.5~\% sur les tickers).
\item Le mauvais lien~: un titre non lié est une transaction manquée, un titre mal lié est une transaction sur le mauvais titre, avec un
rendement qui ressemble à du bruit ou, pire, à un signal.
\item 100~\% sur 737 avec un mot de secteur, 58.0~\% sur 245 sans.
\end{enumerate}
\textbf{Partie IV.}
\begin{enumerate}[resume]
\item \emph{Quel Apex~?} La liaison à date est juste pour 99.5~\% des mentions (précision de 99.5~\%)~; la liaison avec les alias actuels
est juste pour 87.5~\% (précision de 97.5~\% au total, 91.5~\% sur les tickers). Les modèles supervisés par les rendements atteignent un IC de 0.273
(mots TF-IDF, \omterm{def:ml:linear-and-regularised-baselines:logistic}{régression logistique}) et 0.238 (la liste supervisée de mots), contre 0.224 pour le dictionnaire.
\item La liaison~: tout modèle en aval hérite de ses erreurs, et ces erreurs passent inaperçues.
\item Un échantillon de liens stratifié par année et par type de mention, y compris les non liés, à comparer avec les alias à date~;
compter les liens vers des sociétés radiées de la cote ou renommées.
\item L'ajuster et le sélectionner dans des plis en marche avant sur des textes antérieurs, fixer la fenêtre de réaction à l'avance, et le refiltrer selon un
calendrier, en vérifiant la stabilité de la liste.
\item Organiser de longs documents et construire des variables pour un modèle supervisé~; non mesurer le ton.
\item Les jetons par document montent à des milliers, les dictionnaires deviennent plus stables, la structure par sections et les changements par rapport au
dépôt précédent portent de l'information, et une mention peut concerner un pair, un fournisseur ou un client plutôt que le déclarant.
\item Peu d'étiquettes, un nouveau domaine, ou le besoin d'un score transparent~; et des listes construites pour le domaine.
\item Bien identifier le titre.
\end{enumerate}
\end{solution}

\begin{solution}{iq:ml:text-from-bag-of-words-to-embeddings:1}
Comptes de termes pondérés par le logarithme de l'inverse de la part des documents contenant le terme, lignes normalisées. Il réduit les mots courants
vers zéro et ramène les documents de longueurs différentes à une même échelle, de sorte qu'un modèle linéaire pénalisé consacre son budget aux
mots distinctifs.
\iqlookfor{la formule, la raison de l'IDF, et la normalisation selon la longueur des documents.}
\end{solution}

\begin{solution}{iq:ml:text-from-bag-of-words-to-embeddings:2}
Les listes générales classent comme négatifs des mots financiers courants (liability, tax, cost, capital) et manquent leurs sens financiers~; Loughran
et McDonald (2011) ont constaté que près des trois quarts des mots négatifs d'un dictionnaire généraliste largement utilisé ne sont pas
négatifs dans les 10-K, et ont construit des listes financières.
\iqlookfor{mauvaise classification des mots du domaine, avec la référence ou un exemple.}
\end{solution}

\begin{solution}{iq:ml:text-from-bag-of-words-to-embeddings:3}
Sélectionner les mots selon leur association avec les rendements ultérieurs sur les données d'entraînement, les pondérer, noter les documents~; ou ajuster une
\omterm{def:ml:linear-and-regularised-baselines:logistic}{régression logistique} pénalisée sur le TF-IDF. Fuites~: rendements antérieurs à l'horodatage du texte (prendre l'instant où le texte est devenu
disponible), noms d'entreprises tenant lieu de leurs rendements, étiquettes issues de fenêtres qui se chevauchent, et vocabulaire
choisi sur l'échantillon complet.
\iqlookfor{une construction supervisée et au moins deux fuites concrètes avec leurs remèdes.}
\end{solution}

\begin{solution}{iq:ml:text-from-bag-of-words-to-embeddings:4}
Les plongements sont entraînés pour refléter des contextes partagés, et un mot et son contraire occupent les mêmes contextes. Les distances dans
l'espace mesurent le thème et la syntaxe, non la polarité~; utilisées comme variables, elles exigent une couche supervisée ou un \omterm{def:ml:representation-learning:ssl}{ajustement fin} sur la cible.
\iqlookfor{la raison distributionnelle et la conséquence pour le sentiment.}
\end{solution}

\begin{solution}{iq:ml:text-from-bag-of-words-to-embeddings:5}
Repérer les mentions (gazetteer et tagueur), résoudre les candidats à partir d'alias et de tickers valables à la date du texte, lever les ambiguïtés
par le contexte, produire un identifiant permanent assorti d'une confiance, s'abstenir sous un seuil et tout consigner. Modes de défaillance~:
réutilisation de tickers, changements de nom, noms courts partagés, filiales et marques, et un fichier de symboles rafraîchi sans historique.
\iqlookfor{résolution à l'instant de la donnée, désambiguïsation, abstention, et modes de défaillance nommés.}
\end{solution}

\begin{solution}{iq:ml:text-from-bag-of-words-to-embeddings:6}
D'abord l'infrastructure~: horodatages (heure du texte comparée à l'heure de la transaction), changements de liaison des entités, changement de format ou de couverture chez le fournisseur. Ensuite
fuites dans le backtest~: vocabulaire ou étiquettes choisis sur l'échantillon complet, fenêtre de réaction. Puis la dérive~: nouvelles formulations,
encombrement du signal (livre~7, chapitre~28), et enfin la décroissance attendue d'un modèle surajusté.
\iqlookfor{un ordre allant des problèmes de données à la fuite, puis à la décroissance réelle.}
\end{solution}
